\subsection{Datos del par y recurrencia de la acción}\label{sec:par-accion}

Llamemos \(z_{ij}=(r+ih,c+jh)\), \(i,j\in\{0,1\}\), a las cuatro casillas del cuadrado
geométrico. Numeramos las filas raíz por \((r_0,r_1)=(r,r+h)\). En la orientación \(D\)
numeramos las columnas raíz por \((c_0,c_1)=(c,c+h)\) y en la orientación \(A\) por
\((c_0,c_1)=(c+h,c)\). Así se perforan siempre \((r_0,c_0),(r_1,c_1)\). En el orden
de raíces, el bloque de incidencias es \(\left(\begin{smallmatrix}0&X\\Y&0\end{smallmatrix}\right)\).
En coordenadas del tablero esto da la siguiente tabla.
\begin{center}
\begin{tabular}{@{}ccc@{}}
\toprule
Estado \((X,Y)\) & Cruces abiertos para \(D\) & Cruces abiertos para \(A\)\\
\midrule
\((0,0)\)& ninguno & ninguno\\
\((1,0)\)& \(z_{01}\) & \(z_{00}\)\\
\((0,1)\)& \(z_{10}\) & \(z_{11}\)\\
\((1,1)\)& \(z_{01},z_{10}\) & \(z_{00},z_{11}\)\\
\bottomrule
\end{tabular}
\end{center}
Cada fila de la tabla lleva el peso \(w_{o,b}q_{o,(X,Y)}\) del apéndice~\ref{app:boundary}.
Para construir la integral a corte \(b>0\), un tipo ordinario de fila es \((u,\xi)\),
donde \(u\in[0,1]\) y \(\xi=(\xi_0,\xi_1)\) indica las aristas que llegan a las dos columnas
raíz. Escribamos \(e_1(v)=v,\ e_0(v)=1-v\). Las medidas ordinarias de fila y columna son
\[
 d\alpha_R(u,\xi)=du\prod_{k=0}^1 e_{\xi_k}(v_b(u,c_k)),\qquad
 d\alpha_C(v,\zeta)=dv\prod_{k=0}^1 e_{\zeta_k}(v_b(r_k,v)).
\]
A cada lado se añaden dos copias de cada raíz, de masa basal cero. La medida firmada
\(\rho_\sigma\), en el lado \(\sigma=R,C\), asigna masas \(-1,+1\) a sus copias
negativa y positiva. La capacidad \(\lambda\) vale \(1-t\) en las copias positivas y uno
en los demás tipos. Fijamos \(m_\sigma=\lambda\rho_\sigma\),
\(M_\sigma=\alpha_\sigma+s\rho_\sigma\) y \(N_\sigma=\alpha_\sigma+s m_\sigma\).

Sea \(V_{il}\) la incidencia condicional entre un tipo fila \(i\) y un tipo columna \(l\).
Entre tipos ordinarios de posiciones \(u,v\) vale \(v_b(u,v)\). Entre una fila ordinaria y
una copia de la columna raíz \(k\) vale \(\xi_k\). Entre una copia de la fila raíz \(k\) y
una columna ordinaria vale \(\zeta_k\). Entre copias se usa la matriz del estado indicado.
En todos los bloques \(W=W_b^0=p_b/v_b\) se evalúa en las posiciones de ambos tipos.
Con \(q_{il}=\lambda_i\lambda_l,\ \kappa_{1,il}=q_{il}W_{il}\), sustituimos la incidencia por \(V\)
en las ecuaciones de márgenes de este apéndice. Para el determinante, la corrección por
la ida y vuelta en una misma arista se calcula en la sección~\ref{sec:par-prefactor}.

\paragraph{Potenciales y densidades por orden.}
Sean \(\mathcal Z_d=(Z_d,H_d)=[s^d](z,y)\) y \(h_d(i,l)=Z_d(i)+H_d(l)\).
Partimos de \(\mathcal Z_0=0,\ w_0=W\). Para \(d=1,2,3\), el término conocido \(v_d\) de
\(w_d=[s^d]w=Wh_d+v_d\) se obtiene de \(w(1-sqw)=W\exp(z+y)\).
Omitiendo los índices de arista,
\begin{align*}
 v_1/W&=\kappa_1,\\
 v_2/W&=h_1^2/2+2\kappa_1 h_1+2\kappa_1^2,\\
 v_3/W&=h_1h_2+h_1^3/6+\kappa_1(2h_2+2h_1^2)+6\kappa_1^2h_1+5\kappa_1^3 .
\end{align*}
A orden \(d\), los dos márgenes y el calibre determinan \(\mathcal Z_d\) mediante el sistema
\begin{align*}
 Z_d(i)+\int VW H_d(l)\,d\alpha_C(l)
   &=-\int Vv_d\,d\alpha_C-\int Vw_{d-1}\,dm_C ,\\
 H_d(l)+\int VW Z_d(i)\,d\alpha_R(i)
   &=-\int Vv_d\,d\alpha_R-\int Vw_{d-1}\,dm_R ,\\
 \int Z_d\,d\alpha_R&=\int H_d\,d\alpha_C .
\end{align*}
Se resuelve primero el bloque ordinario. La brecha de \(p_b\) da la inversa en el calibre,
y las incógnitas de las copias se recuperan de su ecuación con diagonal uno.
Los lados derechos son compatibles porque las sumas de los dos márgenes respecto
de \(N_R,N_C\) coinciden y estas medidas tienen la misma masa.

\paragraph{Extracción de la acción.}
Escribimos \(\langle\alpha,\mathcal Z_d\rangle=\int Z_d\,d\alpha_R+\int H_d\,d\alpha_C\)
y análogamente para \(m\). Sean \(D_j=[s^j](N_R\otimes N_C)\). En particular,
\(D_0=\alpha_R\otimes\alpha_C,\ D_1=m_R\otimes\alpha_C+\alpha_R\otimes m_C\) y
\(D_2=m_R\otimes m_C\). En cada arista se calculan
\[
 f_1=\tfrac12 qW^2,\qquad f_2=qWw_1+\tfrac13q^2W^3,\qquad
 f_3=q(Ww_2+w_1^2/2)+q^2W^2w_1+q^3W^4/4 .
\]
En la acción dual, extraer los términos de grados uno y dos da
\begin{align*}
 R_1(t)&=-\langle\alpha,\mathcal Z_2\rangle-\langle m,\mathcal Z_1\rangle
          -\int Vf_2\,dD_0-\int Vf_1\,dD_1,\\
 R_2(t)&=-\langle\alpha,\mathcal Z_3\rangle-\langle m,\mathcal Z_2\rangle
          -\int Vf_3\,dD_0-\int Vf_2\,dD_1-\int Vf_1\,dD_2 .
\end{align*}
En efecto, \(f_1,f_2,f_3\) son los tres primeros coeficientes de
\(\sum_{a\ge2}s^{a-2}q^{a-1}w^a/a\), que se integra con signo negativo contra
\(V\,d(N_R\otimes N_C)\). La contribución de los potenciales es
\(-s^{-1}(\int z\,dN_R+\int y\,dN_C)\), módulo los términos afines en \(t\)
que se anulan en el resto de Taylor.
